\section{Dirichlet density of a set of prime ideals}

Throughout, $K$ denotes a number field, $\OK$ its ring of integers, and
$\fp$ ranges over nonzero prime ideals of $\OK$ with absolute norm
$N\fp=[\OK:\fp]$.

\begin{definition}\label{def:dirichlet-density}
  \lean{Chebotarev.HasDirichletDensity}
  \leanok
  A set $S$ of nonzero prime ideals of $\OK$ has \emph{Dirichlet
  density} $\delta\in\R$, written $\delta(S)=\delta$, if
  \[
    \lim_{s\to 1^+}
    \frac{\displaystyle\sum_{\fp\in S} N\fp^{-s}}
         {\displaystyle\sum_{\fp} N\fp^{-s}}
    = \delta,
  \]
  where the denominator sum runs over all nonzero prime ideals of
  $\OK$.
\end{definition}

The denominator is asymptotically $\log\bigl(1/(s-1)\bigr)$ as
$s\downarrow 1$ (Sharifi Prop.~7.1.12, p.~140); the lemmas below break
out the analytic ingredients.

\begin{lemma}\label{lem:prime-zeta-higher-tail-bounded}
  \lean{Chebotarev.primeIdealZetaHigherTail_bounded}
  The higher-power tail of the log-Euler-product
  $\sum_{\fp,k\ge 2}\frac{N\fp^{-2s}}{1-N\fp^{-s}}$ is bounded on a
  right neighbourhood of $s=1$.
\end{lemma}

\begin{lemma}\label{lem:prime-zeta-le-finrank-riemann}
  \lean{Chebotarev.primeIdealZetaSum_le_finrank_smul_riemannPrimeSum}
  \uses{def:dirichlet-density}
  For $s>1$ in a right neighbourhood of $1$,
  \[
    \sum_{\fp} N\fp^{-s} \;\le\; [K:\Q]\sum_{p\text{ rational prime}} p^{-s}.
  \]
\end{lemma}

\begin{proof}
  At most $[K:\Q]$ primes of $\OK$ lie above any rational prime $p$
  (the identity
  $\sum_{\fp\mid p}e(\fp\mid p)f(\fp\mid p)=[K:\Q]$ in mathlib's
  \texttt{Ideal.sum\_ramification\_inertia}), and the norm of such a
  prime is at least $p$, so the bound follows by absolute convergence on
  $\re(s)>1$.
\end{proof}

\begin{lemma}\label{lem:riemann-prime-sum-asymp}
  \lean{Chebotarev.riemannPrimeSum_asymp_log}
  As $s\downarrow 1$,
  \[
    \sum_{p\text{ rational prime}} p^{-s} \sim \log\frac{1}{s-1}.
  \]
\end{lemma}

\begin{proof}
  The case $K=\Q$ of \cref{lem:prime-ideal-sum-log} below: combine the
  Euler product
  $\zeta(s)=\prod_p(1-p^{-s})^{-1}$ with the simple pole of $\zeta$ at
  $s=1$.
\end{proof}

\begin{lemma}\label{lem:logzeta-eq-primesum-bounded}
  \lean{Chebotarev.logDedekindZeta_sub_primeIdealZetaSum_bounded}
  \uses{lem:prime-zeta-higher-tail-bounded}
  Euler-product-log identity:
  $\bigl|\log\zeta_K(s) - \sum_\fp N\fp^{-s}\bigr| \le C$ on a right
  neighbourhood of $1$.
\end{lemma}

\begin{proof}
  \uses{lem:prime-zeta-higher-tail-bounded}
  $\log\zeta_K(s) = \sum_\fp N\fp^{-s} +
  \sum_{\fp,k\ge2} N\fp^{-ks}/k$; the tail is bounded by
  \cref{lem:prime-zeta-higher-tail-bounded}.
\end{proof}

\begin{lemma}\label{lem:logzeta-eq-loginv-bounded}
  \lean{Chebotarev.logDedekindZeta_sub_log_inv_sub_one_bounded}
  Simple-pole identity:
  $\bigl|\log\zeta_K(s) - \log(1/(s-1))\bigr| \le C$ on a right
  neighbourhood of $1$.
\end{lemma}

\begin{proof}
  From the simple pole of $\zeta_K$ at $s=1$ (mathlib's
  \texttt{NumberField.tendsto\_sub\_one\_mul\_dedekindZeta\_nhdsGT}).
\end{proof}

\begin{lemma}\label{lem:prime-zeta-ge-log-minus-bounded}
  \lean{Chebotarev.primeIdealZetaSum_ge_log_minus_bounded}
  \uses{lem:logzeta-eq-primesum-bounded, lem:logzeta-eq-loginv-bounded}
  There is a constant $C$ such that, for $s>1$ in a right neighbourhood
  of $1$,
  \[
    \sum_{\fp} N\fp^{-s} \;\ge\; \log\frac{1}{s-1} - C.
  \]
\end{lemma}

\begin{proof}
  \uses{lem:logzeta-eq-primesum-bounded, lem:logzeta-eq-loginv-bounded}
  \leanok
  Chain the two identities
  \cref{lem:logzeta-eq-primesum-bounded,lem:logzeta-eq-loginv-bounded}:
  $\sum_\fp N\fp^{-s} \ge \log\zeta_K(s) - C_1 \ge \log(1/(s-1)) - C_1 - C_2$.
\end{proof}

\begin{lemma}\label{lem:prime-zeta-le-log-plus-bounded}
  \lean{Chebotarev.primeIdealZetaSum_le_log_plus_bounded}
  \uses{lem:logzeta-eq-primesum-bounded, lem:logzeta-eq-loginv-bounded}
  There is a constant $C'$ such that, for $s>1$ in a right neighbourhood
  of $1$,
  \[
    \sum_{\fp} N\fp^{-s} \;\le\; \log\frac{1}{s-1} + C'.
  \]
\end{lemma}

\begin{proof}
  \uses{lem:logzeta-eq-primesum-bounded, lem:logzeta-eq-loginv-bounded}
  \leanok
  Same two identities
  \cref{lem:logzeta-eq-primesum-bounded,lem:logzeta-eq-loginv-bounded},
  transposed.
\end{proof}

\begin{lemma}\label{lem:tendsto-ratio-one-of-log-pm-bounded}
  \lean{Chebotarev.tendsto_ratio_one_of_log_pm_bounded}
  Let $f : (1,\infty) \to \R$. If $f$ is sandwiched as
  $\log\bigl(1/(s-1)\bigr) - C \le f(s) \le \log\bigl(1/(s-1)\bigr) + C'$
  on a right neighbourhood of $1$, then
  $f(s)/\log\bigl(1/(s-1)\bigr) \to 1$ as $s \downarrow 1$.
\end{lemma}

\begin{proof}
  Since $\log\bigl(1/(s-1)\bigr) \to \infty$ as $s \downarrow 1$, the
  additive constants $C, C'$ wash out under division, and the squeeze
  gives the limit $1$.
\end{proof}

\begin{lemma}\label{lem:prime-ideal-sum-log}
  \lean{Chebotarev.primeIdealZetaSum_univ_tendsto_log}
  \uses{lem:prime-zeta-ge-log-minus-bounded, lem:prime-zeta-le-log-plus-bounded, lem:tendsto-ratio-one-of-log-pm-bounded}
  As $s\downarrow 1$,
  \[
    \sum_{\fp} N\fp^{-s} \sim \log\frac{1}{s-1}.
  \]
\end{lemma}

\begin{proof}
  \uses{lem:prime-zeta-ge-log-minus-bounded, lem:prime-zeta-le-log-plus-bounded, lem:tendsto-ratio-one-of-log-pm-bounded}
  Combine the lower bound \cref{lem:prime-zeta-ge-log-minus-bounded} and
  the upper bound \cref{lem:prime-zeta-le-log-plus-bounded} into a
  two-sided $\log\bigl(1/(s-1)\bigr) \pm O(1)$ sandwich, then apply
  \cref{lem:tendsto-ratio-one-of-log-pm-bounded} to extract the ratio
  limit.
\end{proof}

\subsection*{Density API}

The following routine API lemmas are used by the proofs of the
corollaries in \cref{ch:main}.

\begin{lemma}\label{lem:has-density-finite}
  \lean{Chebotarev.hasDirichletDensity_of_finite}
  \uses{def:dirichlet-density}
  A finite set of prime ideals has Dirichlet density $0$.
\end{lemma}

\begin{proof}
  The numerator $\sum_{\fp\in S}\Norm{\fp}^{-s}$ is bounded as $s \to
  1^+$ (a finite sum of bounded terms), while the denominator
  $\sum_\fp \Norm{\fp}^{-s}\to\infty$ by
  \cref{lem:prime-ideal-sum-log}; the ratio tends to $0$.
\end{proof}

\begin{lemma}\label{lem:density-implies-lower}
  \lean{Chebotarev.HasDirichletDensity.hasLower}
  \uses{def:dirichlet-density}
  If $S$ has Dirichlet density $\delta$, then its lower Dirichlet
  density is also $\delta$.
\end{lemma}

\begin{proof}
  Convergence of the ratio implies the liminf and limsup coincide with
  the limit; specifically the liminf equals $\delta$.
\end{proof}
